\section{Evaluation Setup}
\label{sec:setup}
\label{sec:problem}

Given a passage $x$ and ordered candidates $C=(c_1,\ldots,c_m)$,
admission stores $K\subseteq C$ with $|K|\leq k$. The future question
$q$ and answer $a$ are hidden. When $q$ arrives, the frozen reader sees
only $q$ and $K$; the original passage has left its window. Every policy
receives the same candidates and budget. We measure three successive
outcomes: within-episode ranking, useful-concept retention, and answer
accuracy.

\paragraph{Within-episode ranking.}
Within-episode area under the receiver operating characteristic curve
(AUC) measures how well a policy ranks candidates competing for memory.
It is the fraction of pairs of the useful concept and another candidate
ordered correctly, with half credit for ties, averaged across episodes.
For useful concept $b_i$ in episode $i$, score $s$ and $n$ episodes,
\begin{equation}
\begin{split}
A_{\mathrm{episode}}={}&\frac{1}{n}\sum_{i=1}^{n}\frac{1}{|C_i|-1}\\
&\times\sum_{c\in C_i\setminus\{b_i\}}\Big[\mathbf{1}\{s(b_i)>s(c)\}\\
&\qquad\quad+\tfrac12\,\mathbf{1}\{s(b_i)=s(c)\}\Big].
\end{split}
\label{eq:episodeauc}
\end{equation}
Pooled AUC compares candidates across the entire set and summarizes
score discrimination.

A candidate is stated when a case-insensitive stem match finds it in
the passage and unstated otherwise. Replacing utility labels with
statedness labels gives the \emph{statedness AUC}.

\paragraph{Useful-concept retention.}
Useful-concept retention $B_k$ is the fraction of episodes whose stored
set contains the useful concept:
\begin{equation}
B_k=\frac{1}{n}\sum_{i=1}^{n}\mathbf{1}\{b_i\in K_i\}.
\label{eq:retention}
\end{equation}
A policy can rank it above most distractors and still omit it at a
small budget. Gold-concept selection reads the gold utility label to
rank the labeled useful concept first and appears as a reference row
in the tables.

\paragraph{Answer accuracy.}
Answer accuracy is the fraction of future questions that a fixed,
frozen reader answers correctly from the stored concepts under blind
semantic grading. Accuracy also depends on the reader and the other
stored concepts.

The reader answers greedily, with Qwen3 thinking disabled. A blind
semantic judge receives only the question, gold answer and response,
and accepts a committed correct answer or exact equivalent. A
string-match grader agrees with the judge on 85.1\% of answers, and
SDA exceeds the importance report under both graders
(App.~\ref{app:judge}).

The closed-book reader (question only) and the full-context reader
(whole passage) give reference points without memory and with the
complete original input, respectively.

\paragraph{Comparisons.}
Baselines fall into four groups (App.~\ref{app:baselines}):
\begin{itemize}
\item \emph{Untrained model}: the importance report, the reversed
importance report, a 1--10 importance rating, implied-content wording
(an importance prompt that permits implied content), and
statedness-cued wording (an importance prompt that says the passage
may never mention the useful concept).
\item \emph{Statedness rules}: unstated-first (unstated candidates
first, ties broken at random), unstated-first with report ordering
(each statedness class ordered by the importance report), and the
statedness-adjusted report (the importance report plus one fitted
offset for stated candidates).
\item \emph{Teacher signals}: output likelihood, the workspace
readout, and the combined teacher, each used directly.
\item \emph{Ablated students}: adapters (\S\ref{sec:student}) trained
on labels from one teacher signal.
\end{itemize}
Random selection and published admission criteria (\S\ref{sec:recall})
complete the comparison.

Trained results are means over three training runs; differences between
policies are differences of these means
(Table~\ref{tab:paired}). Per-episode analyses (serving and
useful-concept replacement) use one of these runs, the reference run.
Bold table entries identify SDA.

\paragraph{Evaluation scale.}
\label{sec:breadth}
\label{sec:models}

The main diagnostic comparison covers eight models from three
families: Qwen2.5-7B, OLMo-2-7B, Mistral-7B and Qwen3 at 1.7B, 4B,
8B, 14B and 32B \citep{qwen2025,olmo2,jiang2023mistral,qwen3}.
Outside the checkpoint grid of App.~\ref{app:breadth}, the short
names Qwen2.5-7B, OLMo-2-7B and Mistral-7B denote the instruct
checkpoints Qwen2.5-7B-Instruct, OLMo-2-1124-7B-Instruct
and Mistral-7B-Instruct-v0.3. The extended grid covers 23 base and
post-trained checkpoints from these families and GPT-2
(App.~\ref{app:breadth}).

SDA training covers six models: Qwen3-8B,
Qwen3-14B, Qwen3-32B, Qwen2.5-7B, OLMo-2-7B and Mistral-7B.
Each model supplies its own teacher, adapter and frozen reader.

Training and validation splits contain 175 and 45 episodes,
respectively, from 88 Explicit episodes and 132 episodes in which the
passage cues an unstated useful concept that supplies the answer,
including the 57-episode alternate-generator set.

Development uses Decoupled (68 episodes), Compositional (52), and two
separate 68-episode sets adding confusable or unrelated unstated
distractors to Decoupled. The held-out set and the test set contain
316 and 363 episodes, respectively, while the entity-bridge set
contains 42 episodes. Passage controls and pools reuse the test set's
363 episodes.

Evaluation uses memory budgets $k=1,2,3$. The main comparison uses
the statedness-balanced pool, reporting $k=2$ for all six models.

